\paragraph{Responsabilit�s �ditoriales.} % (�ditoriales, pr�sidence de comit�s scientifiques de congr�s, de programmes...)} 
Je suis �diteur associ� des revues suivantes:
\begin{itemize}
% \item BMC Bioinformatics (jusqu'en 2015), 
% \item Statistical Applications in Genetics and Molecular Biology (SAGMB), 
\item Journal de la Soci�t� Fran�aise de Statistique,
\item Biometrics, 
\item Scandinavian Journal of Statistics (depuis 2018)
\end{itemize}

\paragraph{Relecture d'articles, �valuation de projets.}
\begin{itemize}
 \item Relectures d'articles scientifiques pour des revues de 
 statistique (Electron. J. Stat., Annals Appl. Stat., J. Amer. Stat. Assoc., Stat., Biometrics, Comput. Stat. Data Anal., Stat. Comput.), 
 probabilit�s appliqu�es (J. Appl. Proba., Meth. Comput. Appl. Proba.), 
 bioinformatique (Bioinformatics, BMC Bioinformatics, 
 % IEEE Trans. Comput. Biol. Bioinfo., Stat. Appl. Genet. Mol. Biol., Proteomics
 ) et de biologie (%Tree physiology, 
 PLoS One, Genome Biology, Nuc. Acid Res., Mol. Biol. Evol.).
 \item Relecture d'article pour les �ditions de 2015 � 2019 de l'{\sl International Conference on Artificial Intelligence and Statistics} (AIstat: \url{www.aistats.org/}), les �ditions 2015 � 2018 de {\sl Neural Information and Processing Systems} (NIPS~: \url{nips.cc/}) et l'�dition 2019 de {\sl Inernational Conference on Machine Learning} (ICML~: \url{icml.cc/}). Environ 5 articles par �dition.
 \item Relectures occasionnelles pour des congr�s internationaux (Europ. Conf. Comput. Biol., Recomb, ...).
%  \item Evaluation de 1 ou 2 projets de recherche par an.
\end{itemize}

